\answersfor{ch:bandpipeline}

\begin{solution}{ex:bandpipeline:1}
Weight form. The fraction of the way:
\[ t = \frac{1{,}610 - 1{,}600}{1{,}612 - 1{,}600} = \frac{10}{12} \approx 0.8333. \]
Weights $1 - t \approx 0.1667$ and $t \approx 0.8333$:
\begin{align*}
y &= 0.1667 \times 0.20 + 0.8333 \times 0.26\\
  &= 0.03333 + 0.21667\\
  &= 0.25000.
\end{align*}
Two-point form. Slope:
\[ m = \frac{0.26 - 0.20}{1{,}612 - 1{,}600} = \frac{0.06}{12} = 0.005 \text{ per nm}. \]
Then
\[ y = 0.20 + 0.005 \times (1{,}610 - 1{,}600) = 0.20 + 0.05 = 0.25. \]
Both give $0.25$.
\end{solution}

\begin{solution}{ex:bandpipeline:2}
A pixel is wiped when $u/C' > 0.4$, so it is kept when $u \le 0.4 \times 200 = 80$. The largest number
of invalid bands for a kept pixel is $80$. A pixel with exactly $80$ invalid bands has
$f = 80/200 = 0.4$, and $0.4 > 0.4$ is false, so it is kept and its $80$ holes are filled at stage 7.
Had the code used ``$\ge$'', the same pixel would be wiped: exactly at the threshold, the choice of
comparison decides.
\end{solution}

\begin{solution}{ex:bandpipeline:3}
\code{np.interp} assumes ascending wavelengths and first checks whether the query is at or beyond the
\emph{last listed} wavelength, here $700$. Since $750 \ge 700$, it answers ``beyond the end'' and
returns the last listed value, $0.1$. The correct answer: sort to $(700, 800, 900)$ with values
$(0.1, 0.2, 0.3)$; $750$ lies between $700$ and $800$ with $t = 50/100 = 0.5$, so
$y = 0.5 \times 0.1 + 0.5 \times 0.2 = 0.05 + 0.10 = 0.15$.
\end{solution}

\begin{solution}{ex:bandpipeline:4}
With $\mu = 0.25$, $\sigma = 0.05$:
\begin{align*}
(0.15 - 0.25)/0.05 &= -0.10/0.05 = -2,\\
(0.25 - 0.25)/0.05 &= 0/0.05 = 0,\\
(0.40 - 0.25)/0.05 &= 0.15/0.05 = 3,\\
(0 - 0.25)/0.05 &= -0.25/0.05 = -5.
\end{align*}
The invalid pixel lands at $-5$, far outside the typical range, because the code masks before it
normalises. Denormalise $z = -0.6$:
\[ x = \sigma z + \mu = 0.05 \times (-0.6) + 0.25 = -0.03 + 0.25 = 0.22. \]
Round trip for $0.40$: $N(0.40) = 3$, and $D(3) = 0.05 \times 3 + 0.25 = 0.15 + 0.25 = 0.40$.
Algebraically, $D(N(x)) = \sigma\,\frac{x - \mu}{\sigma} + \mu = x - \mu + \mu = x$ for any $x$.
\end{solution}

\begin{solution}{ex:bandpipeline:5}
Valid positions are $1$ ($0.20$) and $5$ ($0.35$).
\begin{itemize}
\item Position 0: distance $1$ to position $1$, $5$ to position $5$: $0.20$.
\item Position 2: distance $1$ to position $1$, $3$ to position $5$: $0.20$.
\item Position 3: distance $2$ to position $1$ and $2$ to position $5$: a \emph{tie}; the
implementation chooses one of $0.20$ and $0.35$.
\item Position 4: distance $3$ to position $1$, $1$ to position $5$: $0.35$.
\end{itemize}
Result $0.20, 0.20, 0.20, \text{($0.20$ or $0.35$)}, 0.35, 0.35$.

The fill copies a real measurement so that every value the network's first layer sees is physically
plausible. An average of two neighbours could be a spectrum that exists nowhere in the scene, and in a
system that hunts for values that do not fit, an invented value could itself look anomalous.
\end{solution}
